\documentclass[conference,compsoc]{IEEEtran}
% IEEEtran.cls version 1.8b -- US letter paper is the default.
% Do NOT add geometry or other packages that alter margins/paper size.

\ifCLASSOPTIONcompsoc
  \usepackage[nocompress]{cite}
\else
  \usepackage{cite}
\fi

\usepackage{amsmath}
\usepackage{booktabs}
\usepackage{url}

\hyphenation{op-tical net-works semi-conduc-tor}

\begin{document}

\title{Poster: Privacy-Preserving Malware Analysis with On-Prem LLMs
       via Reinforcement Learning Alignment}

\author{%
  \IEEEauthorblockN{Gaspard Baye}
  \IEEEauthorblockA{Chief Technology Officer\\
  Valix AI\\
  baye@valix.ai}
}

\maketitle

% ---------------------------------------------------------------
\begin{abstract}
Security Operations Center (SOC) analysts face a dual challenge: the
volume of AI-enabled threats is growing while data-privacy regulations
prohibit sending sensitive endpoint telemetry to cloud AI APIs.
We present a \emph{preliminary feasibility study} of a three-stage
reinforcement learning (RL) pipeline, combining Supervised Fine-Tuning
(SFT), Proximal Policy Optimization (PPO)~\cite{ppo}, and Group Relative
Policy Optimization (GRPO)~\cite{grpo}, for aligning a compact on-premise
LLM (LLaMA~3.2, 3B~\cite{llama3}) toward analyst-grade malware triage.
Rewards are grounded in MITRE~ATT\&CK~\cite{mitre} terminology and
training data is bootstrapped from a 1-million-sample HuggingFace dataset
of Gemini~\cite{gemini}-distilled reasoning chains, eliminating human annotation.
On a 20-sample pilot benchmark, unaligned open-source baselines already
reach 75--85\% task accuracy, providing a strong starting point for the RL
alignment stages.
To our knowledge, this is the first application of ATT\&CK-grounded
multi-objective reward design to RL alignment of a privacy-preserving
SOC-focused LLM.
Full post-alignment results and qualitative analyses will be presented at
the poster session.
\end{abstract}

% ---------------------------------------------------------------
\section{Introduction}

SOC analysts must rapidly classify malware behaviors from sandbox reports,
map them to MITRE~ATT\&CK tactics, and prioritize incident responses
under an ever-increasing alert volume~\cite{ctibench}.
Frontier LLM APIs (e.g., GPT-4, Gemini) automate this triage effectively
but require sensitive behavioral telemetry to leave the organization's
perimeter, conflicting with compliance frameworks such as SOC~2, HIPAA,
and NIST~800-53.

\textbf{Threat model.}
We consider an adversary capable of intercepting or correlating cloud API
queries to infer the organization's threat posture.
The asset under protection is endpoint and network behavioral telemetry.
An on-premise LLM eliminates this exfiltration surface: inference is
confined inside the security perimeter, no queries cross organizational
boundaries, and the model is inaccessible to external parties.

\textbf{Contributions.}
This work makes three novel contributions:
(1)~the first application of ATT\&CK-grounded multi-objective reward design
to RL alignment for SOC reasoning;
(2)~a \emph{Silver-Standard Distillation} framework that bootstraps
high-quality domain-specific training data from frontier models without
human annotation;
and (3)~an empirical comparison of the accuracy-privacy tradeoff between
cloud frontier models and locally aligned open-source LLMs for malware
triage.

% ---------------------------------------------------------------
\section{Methodology}

Our pipeline follows a ``teach-then-align'' strategy: SFT teaches the model
to reason in structured analyst format, and two successive RL stages sharpen
its accuracy and professional vocabulary.

\textbf{Step~0 -- Silver-Standard Distillation.}
We curate 10,000 raw malware sandbox reports from Hybrid-Analysis
\cite{hybridanalysis} and VirusTotal~\cite{virustotal}.
Gemini~2.0-Flash~\cite{gemini} generates a structured reasoning chain
$c_i$ and threat label $y_i$ per report, which are then semantically
augmented to form a 1-million-sample HuggingFace dataset~$\mathcal{D}$.
Potential bias from Gemini is mitigated by using its outputs only as an
imitation target in Stage~1; verifiable, rule-based rewards in Stages~2--3
steer the model toward correct, analyst-grade behaviors.

\textbf{Stage~1 -- SFT.}
LLaMA~3.2 (3B)~\cite{llama3} is fine-tuned on $\mathcal{D}$
\cite{deepseekmath} to produce structured
\texttt{<thought>}/\texttt{<answer>} output:
\begin{equation}
  \mathcal{L}_\text{SFT}(\theta) = -\!\sum_i \log P_\theta(y_i,\,c_i \mid x_i).
\end{equation}
This instills reasoning structure the base model lacks entirely.

\textbf{Stage~2 -- PPO~\cite{ppo}.}
A composite reward is maximized while a frozen reference policy
$\pi_\text{ref}$ (KL anchor) prevents reward exploitation.

\textbf{Stage~3 -- GRPO~\cite{grpo}.}
$G{=}8$ candidate responses are sampled per report; within-group advantage
$A_k{=}(r_k{-}\mu(r))/\sigma(r)$ drives the policy toward the best
response without requiring a separate critic network, making it
memory-efficient for local hardware.

\textbf{Reward architecture.}
The composite reward mirrors three criteria a SOC analyst applies when
evaluating a malware report:
\begin{equation}
  R(O\!\mid\! Q)\!=\!
    \underbrace{\lambda_\text{fmt}R_\text{fmt}}_{\text{format}}
  +\underbrace{\lambda_\text{acc}R_\text{acc}}_{\text{accuracy}}
  +\underbrace{\lambda_\text{prec}R_\text{prec}}_{\text{ATT\&CK terms}},
\end{equation}
where $R_\text{fmt}\!\in\![0,0.5]$ rewards \texttt{<thought>}/\texttt{<answer>}
adherence; $R_\text{acc}\!\in\!\{0,1\}$ rewards exact ground-truth label
match; and $R_\text{prec}$ scores the density of high-value MITRE~ATT\&CK
\cite{mitre} terms (e.g., ``process injection'', ``C2 beacon'',
``persistence'') in the model's response.

% ---------------------------------------------------------------
\section{Preliminary Results}

Table~\ref{tab:results} reports pilot performance before any RL alignment.
Results should be interpreted as a feasibility signal on a small sample
($N{=}20$); statistical evaluation with confidence intervals will accompany
the full study.
LLaMA-3 (8B)~\cite{llama3} reaches 85\% accuracy with highly concise
208-word responses, demonstrating strong label-extraction capability.
Mistral (7B)~\cite{mistral} achieves 75\% but exhibits higher ATT\&CK term
density (0.091 vs.\ 0.082), suggesting richer security vocabulary at the
cost of accuracy -- a tradeoff the joint reward architecture is designed to
resolve.
Critically, \emph{no baseline produces structured reasoning}
(Reasoning Density\,$=\!0$); all answers lack explicit step-by-step
justification.
The most common failure mode observed is \emph{partial label extraction}:
the correct malware family is identified but secondary behaviors
(e.g., associated persistence mechanism) are missed.
The \texttt{<thought>} chain and GRPO reward directly target this gap.

\begin{table}[!t]
  \renewcommand{\arraystretch}{1.15}
  \caption{Pilot Evaluation ($N{=}20$). Gemini is a cloud reference only;
           its score reflects ideal prompting on a small sample.
           Aligned model results will be presented at the poster session.}
  \label{tab:results}
  \centering
  \begin{tabular}{lcccc}
    \toprule
    \textbf{Model} & \textbf{Acc.} & \textbf{Prec.} &
    \textbf{Words} & \textbf{Reasoning} \\
    \midrule
    Gemini~2.0~\cite{gemini} (cloud ref.) & 98\% & 0.108 & 505 & 0.00 \\
    LLaMA-3 8B~\cite{llama3} (baseline)   & 85\% & 0.082 & 208 & 0.00 \\
    Mistral 7B~\cite{mistral} (baseline)  & 75\% & 0.091 & 282 & 0.00 \\
    \midrule
    SFT / PPO / GRPO Expert & \multicolumn{4}{c}{\emph{poster session}} \\
    \bottomrule
  \end{tabular}
\end{table}

% ---------------------------------------------------------------
\section{Conclusion and Reproducibility}

This work establishes a three-stage RL alignment methodology for
privacy-preserving SOC malware triage, demonstrating that compact
open-source LLMs are viable starting points before alignment.
We make no claim of parity with frontier models; our goal is to define
the methodology and characterize the baseline trajectory.
Full statistical evaluation, ablation studies, and robustness testing
against adversarial prompt injection will be presented at the poster session.
\textbf{Setup:} LLaMA~3.2~(3B)~\cite{llama3}, HuggingFace
TRL~\cite{trl}, 1$\times$NVIDIA H100; SFT $\text{lr}{=}2{\times}10^{-4}$,
PPO/GRPO $\text{lr}{=}5{\times}10^{-6}$, $G{=}8$.
Training prompts, reward source code, and evaluation scripts will be
released upon acceptance to support community reproducibility.

\ifCLASSOPTIONcompsoc
  \section*{Acknowledgments}
\else
  \section*{Acknowledgment}
\fi
The author thanks HuggingFace TRL~\cite{trl}, the LLaMA
project~\cite{llama3}, Hybrid-Analysis~\cite{hybridanalysis}, and
VirusTotal~\cite{virustotal}.

% ---------------------------------------------------------------
\begin{thebibliography}{11}

\bibitem{ppo}
J.~Schulman \emph{et al.}, ``Proximal Policy Optimization Algorithms,''
\emph{arXiv:1707.06347}, 2017.

\bibitem{grpo}
Z.~Shao \emph{et al.}, ``DeepSeekMath: Pushing the Limits of Mathematical
Reasoning in Open Language Models,'' \emph{arXiv:2402.03300}, 2024.

\bibitem{deepseekmath}
D.~Guo \emph{et al.}, ``DeepSeek-R1: Incentivizing Reasoning Capability
in LLMs via Reinforcement Learning,'' \emph{arXiv:2501.12948}, 2025.

\bibitem{trl}
L.~von Werra \emph{et al.}, ``TRL: Transformer Reinforcement Learning,''
\url{https://github.com/huggingface/trl}, 2020.

\bibitem{llama3}
A.~Dubey \emph{et al.}, ``The Llama 3 Herd of Models,''
\emph{arXiv:2407.21783}, 2024.

\bibitem{mistral}
A.~Q.~Jiang \emph{et al.}, ``Mistral 7B,'' \emph{arXiv:2310.06825}, 2023.

\bibitem{gemini}
Google DeepMind, ``Gemini 2.0,''
\url{https://deepmind.google/technologies/gemini/}, 2024.

\bibitem{mitre}
MITRE Corporation, ``ATT\&CK,'' \url{https://attack.mitre.org}, 2015.

\bibitem{ctibench}
M.~Ameri \emph{et al.}, ``CTIBench,'' \emph{arXiv:2406.07521}, 2024.

\bibitem{hybridanalysis}
CrowdStrike, ``Hybrid Analysis,'' \url{https://www.hybrid-analysis.com}

\bibitem{virustotal}
Google, ``VirusTotal,'' \url{https://www.virustotal.com}

\end{thebibliography}

\end{document}
